\section*{Scope, reading plan, and notation}
\addcontentsline{toc}{section}{Scope, reading plan, and notation}
This guide follows your topic list. The repeated expectation item is combined into one topic; the seven resulting sections have ten original questions each. These questions model methods from Homeworks 4--5, rather than predicting the actual exam. The earlier whole-course study guide serves as a reading map; the textbook and homework PDFs supply the mathematical sources checked for this edition.

\textbf{On an exam.} Start with \texttt{EXAM ASSISTANT} or \texttt{GENERAL SOLVERS}. Select the overall task and enter the current question's values, formula, support, and event. Homework presets are practice examples, not exam predictions.

\textbf{Study sequence.} Read a concept and its example, cover the checkpoint answer, and solve the checkpoint. Then attempt the corresponding ten-question practice set without consulting the solutions. Use the homework keys to check whether you can reproduce the assigned methods. Finally work a mixed set: A2, A6, B4, C6, D2, D6, E4, F3, F8, G3.

\begin{tabularx}{\linewidth}{@{}p{0.10\linewidth}Xp{0.31\linewidth}@{}}
\toprule Set & Skill & Homework connection\\\midrule
A & Joint, marginal, conditional distributions; support; probabilities & HW4 common 1--3; STA431 4--5\\
B & Expectation, functions of variables, linearity & HW5 common 1--2\\
C & Raw and central moments; variance & HW5 common 1--2\\
D & MGFs, derivatives, transformations, uniqueness & HW5 common 3--4\\
E & Product moments, covariance, correlation & HW5 STA431 8\\
F & One and two linear combinations & HW5 common 5--6; Freund Ch4 \S7\\
G & Conditional mean and variance & HW5 common 7; STA431 8\\\bottomrule
\end{tabularx}

\subsection*{Notation that prevents common errors}
$p_{X,Y}$ denotes a joint PMF, $f_{X,Y}$ a joint PDF, and $F_{X,Y}$ a joint CDF. Freund often uses $f$ for either a PMF or a PDF and $g,h$ for marginals; the meaning depends on the support and whether you sum or integrate. A support is the set on which probability mass or density can be positive. Every density is zero outside its stated support unless otherwise specified.

$\mu_X=\E X$, $\sigma_X^2=\Var(X)$, $\mu'_r=\E(X^r)$, and $\mu_r=\E[(X-\mu_X)^r]$. The quantity $\E(XY)$ is not generally $\E X\E Y$. In conditional formulas the value after the vertical bar is held fixed.

\begin{insight}[A repeatable solution method]
\textbf{1. Model:} define the variables and translate words into an event or payoff.\quad
\textbf{2. Support:} list the pairs or sketch the region.\quad
\textbf{3. Normalize:} check nonnegativity and total mass one.\quad
\textbf{4. Choose:} marginal, conditional, probability, moment, or MGF.\quad
\textbf{5. Set up:} write the full sum/integral with limits before calculating.\quad
\textbf{6. Check:} probability in $[0,1]$, variance nonnegative, sensible units, conditionals normalized.
\end{insight}

\textbf{Handwritten sheet.} Copy the separate landscape master using its three-column layout on each side. The back prioritizes calculator workflows, with native commands and the STAM2 program entry point. The explanations and solutions stay here for learning. US Letter landscape is $11\times8.5$ inches; a 0.3 mm pen helps with small lettering, but handwriting size determines fit. The master uses approximately 8-point body text; use a printout at actual size as your spacing guide.
